\section*{Chapter \ref{ch:m1:reading-market-data} --- Reading Market Data}

\begin{solution}{exo:m1:reading-market-data:1}
Order 1001 disappears (300 executed of 300); order 1002 goes from 200 to 100.
\omterm{def:m1:reading-market-data:levels}{Level 2}: bid 99.99 for 100 in 1 order, 99.98 for 100; asks unchanged. \omterm{def:m1:reading-market-data:levels}{Level
1}: 99.99 $\times$ 100 bid, 100.01 $\times$ 500 offered, last 99.99 $\times$ 300. The
trade printed at 99.99, the resting order's price, which the execution
message does not carry; the resting order was a bid, so the aggressor was a
seller.
\end{solution}

\begin{solution}{exo:m1:reading-market-data:2}
\texttt{0x000F41DC} $= 999\,900$: \$99.9900. Read little-endian the same bytes
are $3\,695\,251\,200$: \$369\,525.12. Byte order errors are usually that
obvious; an error in the number of implied decimals (99.99 read as 9\,999) is
not, in a feed that mixes instruments.
\end{solution}

\begin{solution}{exo:m1:reading-market-data:3}
Each message costs its length plus 2 bytes of framing: $1\,102 \times 38 + 311
\times 33 + 114 \times 25 + 473 \times 21 + 20 \times 46 = 41\,876 + 10\,263 + 2\,850 +
9\,933 + 920 = 65\,842$ bytes.
\end{solution}

\begin{solution}{exo:m1:reading-market-data:4}
$2^{48}$ nanoseconds is 78.2 hours: more than a day, with room for sessions
that cross midnight. Four bytes hold only 4.3 seconds of nanoseconds. Against
eight bytes, six save two bytes on every message, a few percent of a feed of
tens of millions of messages a day, at the cost of an unaligned field that
every decoder must assemble by hand.
\end{solution}

\begin{solution}{exo:m1:reading-market-data:5}
No. The difference is 40 microseconds and each clock may be off by up to 100,
in either direction: the true order is undetermined. To say which moved
first one needs one clock for both: \omterm{def:m1:reading-market-data:ts}{receive timestamps} from a single capture
point with known, stable latencies to each venue, or venues synchronised an
order of magnitude more tightly than the gap being measured.
\end{solution}

\begin{solution}{exo:m1:reading-market-data:6}
Line arbitration: it takes 5\,001 from whichever line delivers first, 5\,002
from A, 5\,003 from B, 5\,004 from either, and the book is complete. If both
lines skip 5\,003 the message is lost: request a retransmission, or wait for
the next snapshot and rebuild, buffering later messages meanwhile. Until
then the book for the affected instruments is marked stale; the strategy
must stop quoting from it, since it no longer knows its own queue positions
or even the best price, and may only reduce risk using a slower, independent
price source.
\end{solution}

\begin{solution}{exo:m1:reading-market-data:7}
About 1\% at 100 microseconds and 11\% at 1 millisecond. With every second
message dropped the gaps between consecutive messages double, so a given
jitter overtakes fewer neighbours and the percentages fall to roughly half
at small jitters. Reordering depends on jitter \emph{relative to the
inter-message time}: it is worst in bursts, which is when order matters
most.
\end{solution}

\begin{solution}{exo:m1:reading-market-data:8}
(i) Bad ticks: decimal errors and cancelled trades that the vendor's bars
still contain produce ``falls'' of 50\% or more that recover by the next bar;
no one could have bought there. (ii) Conditions: late and out-of-sequence
reports, or prints from another day's price level, entering the low. (iii)
\omterm{def:m1:reading-market-data:refdata}{Symbology}: bars keyed by ticker join two companies when a ticker is reused,
and unadjusted splits look like a 50\% fall at the open. One may add
survivorship in the ticker list, and the fact that a bar's low says nothing
about the size available at that price.
\end{solution}

\begin{solution}{pb:m1:reading-market-data:1}
\textbf{1.} 22\,650 shares at 98.40.
\textbf{2.} 10.01.
\textbf{3.} $98.40/100.05 - 1 = -165$ basis points.
\textbf{4.} Explode: two consecutive returns of $-90\%$ and $+900\%$ dominate
any estimate for months.
\textbf{5.} A decimal or keying error: a price ten times too small between
two normal prints a second apart. A filter on distance from a rolling median
or from the prevailing quote catches it; confirm against the quote at
09:30:30 (no bid near 10 existed) and against the exchange's later
cancellation.
\textbf{6.} Its price is normal. Only its condition, or a separate
cancellation message referring to its \omterm{def:m1:reading-market-data:seq}{sequence number}, removes it: step 1 of
\cref{met:m1:reading-market-data:clean}.
\textbf{7.} No. It is volume of the day, and belongs in a daily total, but it
was not executed during the order's life and its price reflects 09:24; a
benchmark for 09:30 to 09:32 must use the execution time, not the report
time.
\textbf{8.} In: it is a real trade in the interval at a market price.
(Whether \omterm{def:m1:us-equity-market-structure:lots}{odd lots} update the consolidated last price is a rule of the tape;
for a volume-weighted average they count.)
\textbf{9.} Out, for this client: the auction closed before the order existed
and could not have been joined. It would be in for an order entered before
the open.
\textbf{10.} Three excluded (prints 6, 9 and 11); 16\,950 shares; 100.045.
\textbf{11.} 4\,950 shares at 100.106.
\textbf{12.} Low 100.02 (the auction) or 100.04 in continuous trading; high
100.15.
\textbf{13.} Against 100.045: $+4.5$ basis points (it paid more). Against
100.106: $-1.6$ (it paid less).
\textbf{14.} The second. A benchmark should contain the trades the order
could have taken part in; the auction is 71\% of the first figure's volume
and was over before the order arrived.
\textbf{15.} As a cancel or break message later in the day, or in the next
day's corrections file, referring to the original trade by its sequence or
match number. The loader must apply it to the stored trade, keep both
records, and recompute anything derived.
\textbf{16.} Because there is always a set of exclusions that makes a given
execution look good. A rule chosen after seeing the result is not a
benchmark.
\textbf{17.} The exact list of conditions excluded, the treatment of
corrections and late reports, the timestamp used for bar assignment, and
whether \omterm{def:m1:us-equity-market-structure:lots}{odd lots} and off-exchange trades are in.
\textbf{18.} Capture: nothing, everything is kept. Normalise: conditions
mapped to explicit flags, corrections linked to originals. Book building:
not involved for trades, but the prevailing quote it produces is what
detects print 6. Tables: the purpose-specific inclusion rule and the outlier
flag.
\textbf{19.} \textbf{Three prints excluded; 100.045 (100.106 for continuous
trading only).}
\textbf{20.} A message saying that someone reported a trade, with a price, a
size, a time and a set of qualifications that decide what it means.
\end{solution}

\begin{solution}{iq:m1:reading-market-data:1}
\omterm{def:m1:reading-market-data:levels}{Level 1}: best bid and offer with sizes, and last trade. \omterm{def:m1:reading-market-data:levels}{Level 2}: total size
at each price, by level. \omterm{def:m1:reading-market-data:levels}{Level 3}: every order and every event on it. Each is
computable from the next; only \omterm{def:m1:reading-market-data:levels}{level 3} gives queue position, order
lifetimes and the identity of what was executed.
\iqlookfor{the one-way relation and what only \omterm{def:m1:reading-market-data:levels}{level 3} provides.}
\end{solution}

\begin{solution}{iq:m1:reading-market-data:2}
A hash map from order reference to the order (side, price, remaining size)
and, per instrument and side, a sorted structure from price to aggregate
size (and, for queue position, the orders at that price in time order). Add:
insert in both. Execute or cancel: look up by reference, reduce both, remove
when empty. Delete: the same with the full size. Best prices are the ends of
the sorted structure. Production books replace trees by arrays indexed by
tick around the touch, and pre-allocate.
\iqlookfor{lookup by reference, because executions carry no price.}
\end{solution}

\begin{solution}{iq:m1:reading-market-data:3}
First check the other line. If both missed it: mark the affected books
stale, tell the strategies, buffer what follows, request retransmission or
take the next snapshot, rebuild, replay the buffer, clear the flag. Log the
gap with times. Trading on a stale book is forbidden; pulling quotes is the
default.
\iqlookfor{the stale flag reaching the strategy.}
\end{solution}

\begin{solution}{iq:m1:reading-market-data:4}
The \omterm{def:m1:reading-market-data:ts}{receive timestamp}, at the point where the strategy would have seen the
data, plus the strategy's own latency to act. \omterm{def:m1:reading-market-data:ts}{Exchange timestamps} give the
backtest information before it could have arrived: it trades on quotes that
were already gone, and cross-venue signals look profitable because one
venue's clock runs ahead of the other's. Exchange time remains right for
ordering events within one venue.
\iqlookfor{look-ahead through the clock.}
\end{solution}

\begin{solution}{iq:m1:reading-market-data:5}
Identifiers: what keys the file, and are they stable through time. Time:
which clock, which zone, execution or report time. Conditions: which are
present, which were removed. Corrections: applied or not. Units: currency,
price scale, adjusted or raw for splits. Duplicates and gaps: per day counts
against an independent total. Outliers against quotes. Coverage: which
venues, on- and off-exchange. Then ten minutes looking at one day by eye.
\iqlookfor{conditions, corrections, identifiers, clock.}
\end{solution}

\begin{solution}{iq:m1:reading-market-data:6}
Store the raw capture, unmodified, with hardware \omterm{def:m1:reading-market-data:ts}{receive timestamps},
compressed and checksummed, by venue, line and day: it is the only ground
truth. Version the decoder and the \omterm{def:m1:reading-market-data:refdata}{reference data}; record with every derived
table the capture files, code version and parameters that produced it.
Derived data (books, bars, tables) are caches that can be deleted and
rebuilt. \omterm{def:m1:reading-market-data:refdata}{Reference data} is kept as of each date, never overwritten. Test
reproducibility by rebuilding a random old day every week and comparing
hashes.
\iqlookfor{raw capture as ground truth; point-in-time \omterm{def:m1:reading-market-data:refdata}{reference data};
provenance on every table.}
\end{solution}
